\chapter{The synthetic generation system, reconstructed from its output}
\label{ch:rules}

\begin{saybox}[title=What to say at the start of this section]
Be explicit that this chapter reconstructs the generator's \emph{behaviour}
from the data it produced, because the generator's code is not in this
repository. Then show the rule inventory --- it is the single most
informative table in the report, and it makes the rest of the argument
almost self-evident.
\end{saybox}

\section{What a \texttt{source\_rule\_id} actually is}

Each row of the V7 corpus carries a \texttt{source\_rule\_id}. Reading the
actual values shows immediately that a rule identifier is \emph{not} a
single publication and \emph{not} a single physical relationship. It is a
pipe-delimited composite of source components. Examples taken verbatim
from the soft specialist:

\begin{center}\small
\begin{tabular}{ll}
\toprule
\texttt{source\_rule\_id} & rows \\
\midrule
\texttt{V7|PH\_CHEESE} & 5\,654 \\
\texttt{V7|PLANT\_REVIEW} & 2\,541 \\
\texttt{V7|RICOTTA\_MAP} & 283 \\
\texttt{V7|NATAMYCIN\_REVIEW} & 152 \\
\texttt{V7|RICOTTA\_MAP|PLANT\_REVIEW} & 145 \\
\texttt{V7|HMMC\_NAT|NATAMYCIN\_REVIEW} & 17 \\
\texttt{V7|RICOTTA\_MAP|NATAMYCIN\_REVIEW} & 8 \\
\bottomrule
\end{tabular}
\end{center}

Three structural properties follow directly, and all three matter for the
evaluation design:

\begin{enumerate}[leftmargin=1.5em, itemsep=2pt]
\item \textbf{Rules share components.} \texttt{RICOTTA\_MAP} appears alone,
  with \texttt{PLANT\_REVIEW}, and with \texttt{NATAMYCIN\_REVIEW}.
  \texttt{NATAMYCIN\_REVIEW} appears in three different composites. Rules
  are therefore not mutually independent even as labels --- they are
  intersections drawn from a smaller set of underlying source components.
\item \textbf{Composite rules are small.} The single-component rules are
  large ($5\,654$ and $2\,541$ rows); the three-way composites are tiny
  ($17$ and $8$ rows). The composites appear to represent the specific
  intersections where two sources both constrain the same formulation.
\item \textbf{The distribution is extremely unbalanced.} The largest rule
  in soft/general is $707$ times the size of the smallest.
\end{enumerate}

\begin{limitbox}[title=What cannot be said]
Because the generator is absent, this report cannot state the mathematical
form of any rule's target equation, its parameter distributions, its noise
model, or whether two rules that share a component also share calibration
constants. The composite naming makes shared structure \emph{likely}, but
that is an inference from identifiers, not a verified property of the
generating code. \Cref{ch:limits} treats this as a live limitation of the
revised evaluation, not a closed question.
\end{limitbox}

\section{Complete rule inventory for all six specialists}

\Cref{tab:tab_rule_inventory} enumerates every generation rule in the
production corpus with the statistics that can be computed from the data:
row count, number of distinct contexts, number of distinct cheese products
and indicator types, the storage-temperature range spanned, and the mean
and standard deviation of the target. \Cref{tab:tab_rule_products} gives
the products and treatment ingredients associated with each rule.

\tbllong{tab_rule_inventory}{Generation-rule inventory for all six
specialists. Every column is computed from the V7 specialist files;
the target equations behind these rules are not recoverable here.}

\tbllong{tab_rule_products}{Cheese products and primary treatment ingredients
associated with each generation rule, as they appear in the data.}

\fig{fig06_rule_size_distribution}{0.96}{%
  Rows per generation rule, logarithmic scale, for each of the six
  specialists. The imbalance is a property of the corpus, not of the
  evaluation, and it is the reason the mean and median of per-fold scores
  diverge so sharply in \Cref{ch:loro}.}

\section{The dependency between rows generated by one rule}

Consider two rows drawn from the same rule. They differ in their feature
values --- a different storage temperature, a different concentration, a
different indicator threshold --- so they are distinct points in feature
space and no row-level duplication exists. What they share is the
\emph{mapping} from those features to the target.

Write the generator for rule $g$ as a function $\phi_g$ with parameters
$\theta_g$ fixed by that rule's source constraints, applied to a sampled
input $\mathbf{x}$ and perturbed by noise:
\begin{equation}
y \;=\; \phi_g\!\left(\mathbf{x}; \theta_g\right) + \varepsilon,
\qquad \mathbf{x} \sim \mathcal{X}_g, \qquad \varepsilon \sim \mathcal{E}_g .
\label{eq:generator}
\end{equation}
The form of $\phi_g$ is unknown to this report. But regardless of its form,
the following holds: if the training set contains many draws from
$\mathcal{X}_g$ and the test set contains further draws from the same
$\mathcal{X}_g$ under the same $\phi_g$ and $\theta_g$, then a sufficiently
flexible learner can recover $\phi_g$ to within $\mathcal{E}_g$ and will
score on the test rows as well as the noise floor allows.

The learning problem in that situation is \emph{interpolation within a
known mechanism}. It is a real learning problem --- the model still has to
recover $\phi_g$ from finitely many samples --- but it is a fundamentally
easier one than predicting a mechanism never seen, and it is not the
problem a deployed shelf-life predictor faces when it meets a new
formulation.

\fig{fig18_context_split_diagram}{0.95}{%
  The original evaluation protocol in schematic form. Grouping by
  \texttt{context\_id} guarantees that no single context contributes rows
  to two partitions. It places no constraint whatever on the generation
  rule, so rows from rule $g$ reach training, validation and test
  simultaneously.}

\section{Contexts are not mechanisms}

The protocol's grouping variable, \texttt{context\_id}, identifies a
formulation-plus-condition scenario; one context contributes several rows
(different indicators, different time points). The V7 soft/general
specialist has $8\,800$ rows spread over its contexts, and each generation
rule spans many contexts.

That nesting is the entire problem. Holding out a context removes one
sample of a mechanism; it does not remove the mechanism. The next chapter
quantifies exactly how much mechanism-level information survives the
original split, and the answer turns out to be: all of it.
